\chapter*{Prefacio}
\addcontentsline{toc}{chapter}{Prefacio}
\markboth{PREFACIO}{PREFACIO}

El cálculo de Malliavin es un cálculo diferencial en el espacio de trayectorias del movimiento
browniano. Nació en 1976, cuando Paul Malliavin dio una prueba probabilística del teorema
de hipoelipticidad de Hörmander, y ha crecido hasta convertirse en una de las herramientas
centrales del análisis estocástico: permite demostrar que las leyes de soluciones de ecuaciones
estocásticas tienen densidad regular, calcular sensibilidades de precios de derivados financieros
sin diferenciar funciones de pago discontinuas, obtener teoremas centrales del límite cuantitativos
(el método de Malliavin--Stein) y, más recientemente, construir y renormalizar soluciones de
ecuaciones en derivadas parciales estocásticas singulares.

Esta guía recorre ese camino completo. La Parte~I establece los fundamentos: martingalas,
movimiento browniano, integral de Itô y ecuaciones diferenciales estocásticas. La Parte~II
desarrolla el núcleo de la teoría en el marco de los procesos gaussianos isonormales: caos de
Wiener, la derivada $D$, la divergencia $\delta$, el operador de Ornstein--Uhlenbeck $L$ y las desigualdades
de Meyer. La Parte~III contiene la aplicación fundacional: criterios de existencia y regularidad
de densidades y el teorema de Hörmander. Las Partes~IV y~V desarrollan con detalle dos
grandes campos de aplicación, las finanzas matemáticas y las EDPEs. La Parte~VI presenta
extensiones (procesos gaussianos no markovianos y espacios de Poisson) y la Parte~VII describe
el estado del arte.

\paragraph{Convenciones.} Las demostraciones son completas para los resultados estructurales; cuando
un argumento es muy técnico se da un esbozo detallado y una referencia precisa. Los recuadros
\emph{Intuición} pueden saltarse en una primera lectura, pero recomendamos no hacerlo. Las constantes
$C$ pueden cambiar de línea a línea. Trabajamos sobre un espacio de probabilidad completo
$(\Omega,\F,\P)$.

\paragraph{Lecturas de referencia.} Las fuentes principales son Nualart~\cite{Nualart06}, Nualart--Nualart~\cite{NualartNualart18},
Nourdin--Peccati~\cite{NourdinPeccati12}, Karatzas--Shreve~\cite{KaratzasShreve91}, Revuz--Yor~\cite{RevuzYor99}, Da Prato--Zabczyk~\cite{DaPratoZabczyk14}, Walsh~\cite{Walsh86}
y Hairer~\cite{Hairer11}.
